% ── Úvod: typologie evropských modelů sociálního dialogu ─────────────────────

\Acfgen{KV} v~Evropě nabývá několika forem podle role státu a~sociálních partnerů; akademická
literatura je obvykle člení do~\emph{čtyř hlavních modelů}, které se liší úrovní
vyjednávání, mírou koordinace mezi odvětvími a~mechanismem rozšiřování \acs{KS}
na~nepokryté firmy:
\emph{nordický korporativismus} (Dánsko, Švédsko, Finsko) postavený na~vysoké
odborové organizovanosti a~systému Ghent; \emph{germánský konsensuální model} (Německo,
Rakousko, Nizozemsko) opřený o~odvětvové smlouvy a~spolurozhodování; \emph{románský
státně regulovaný model} (Francie, Itálie, Španělsko) s~dominantní rolí státu
jako ručitele rozšiřování \acs{KS}; a~\emph{liberální model} (Velká Británie, Irsko,
částečně postkomunistické země včetně~\acgen{ČR}) s~převahou podnikové úrovně
a~slabou koordinací.~\cite{Hala_SystemySocialnihoDialogu2013}

V~následujících čtyřech podkapitolách jsou popsány vybrané modely, které jsou
relevantní pro~komparaci s~\acgen{ČR}: dánský (referenční pro debatu o~flexicurity),
rakouský (geograficky a~strukturálně nejbližší kandidát pro~přenos mechanismů
\acs{KV} \textit{erga omnes}), německý (klíčový obchodní partner \acgen{ČR},
průmyslová ekonomika srovnatelné struktury) a~konečně polský se~slovenským
(nejbližší postkomunistické srovnání). Každá podkapitola zahrnuje
právní a~organizační uspořádání, aktuální stav pokrytí \acs{KS} a~odborové organizovanosti
(viz obr.~\ref{fig:eu_pokryti_kv_mapa} a~\ref{fig:eu_hustota_mapa}), klíčové
rysy modelu a~implikace pro~\acgen{ČR}.

Typologie zároveň zachycuje základ pravidel a~vyjednávacích mechanismů tzv.~dvourychlostní \aclgen{EU}. Nejde pouze o~rozdíl mezi bohatšími a~chudšími zeměmi, nýbrž o~rozdíl mezi ekonomikami, které dokázaly spojit otevřený trh s~kolektivní koordinací mezd a~pracovních standardů, a~ekonomikami, v~nichž otevřenost trhu funguje převážně jako tlak na~nákladovou konkurenceschopnost.\cite{EC_WhitePaper_FutureOfEurope_2017}~\cite{Brixova_MICT_CEE_2026}\par
